% ============================================================
%  CV — Maciej Górski   (wersja elegancka / modern)
%  Kompilacja: XeLaTeX lub LuaLaTeX (zalecane)
%  Overleaf:   Menu → Compiler → XeLaTeX
%
%  UWAGA: Jeśli nie masz fontu Inter, użyj Noto Sans:
%    \setmainfont{Noto Sans}
%  (Noto Sans jest pre-instalowany na Overleaf)
% ============================================================
\documentclass[11pt,a4paper,sans]{moderncv}

% ---------- styl ----------
\moderncvstyle{banking}
\moderncvcolor{gold}
\usepackage[scale=0.82]{geometry}
\setlength{\hintscolumnwidth}{2.6cm}

% ---------- pakiety ----------
\usepackage[polish]{babel}
\usepackage{fontspec}
\setmainfont{Inter}[
  Extension      = .ttf,
  UprightFont    = *-Regular,
  BoldFont       = *-Bold,
  ItalicFont     = *-Italic,
  BoldItalicFont = *-BoldItalic,
]
\usepackage{amssymb}
\usepackage{fontawesome5}
\usepackage{xcolor}
\usepackage{hyperref}
\usepackage{tcolorbox}
\tcbuselibrary{skins}

% ---------- kolory ----------
\definecolor{sidebar}{HTML}{0a0f1c}
\definecolor{sidebarText}{HTML}{e2e8f0}
\definecolor{goldAccent}{HTML}{e5b05c}
\definecolor{goldLight}{HTML}{f3d9a0}
\definecolor{goldDark}{HTML}{c9943a}
\definecolor{mutedText}{HTML}{94a3b8}
\definecolor{softText}{HTML}{cbd5e1}
\definecolor{accentBlue}{HTML}{60a5fa}

% ---------- komenda do linii dekoracyjnej ----------
\newcommand{\goldline}{{\color{goldAccent!40}\rule{1.8em}{0.4pt}}}
\newcommand{\goldlineFull}{{\color{goldAccent!25}\rule{\linewidth}{0.3pt}}}

% ---------- styl sekcji ----------
\renewcommand*{\sectionfont}{\Large\bfseries\color{sidebar}}
\renewcommand*{\subsectionfont}{\large\bfseries\color{black!85}}
\renewcommand*{\hintfont}{\small\color{goldDark}}

% ---------- dokument ----------
\begin{document}

% ======================== PASEK BOCZNY ========================
\begin{tcolorbox}[
  colback      = sidebar,
  coltext      = sidebarText,
  boxrule      = 0pt,
  leftrule     = 0pt,
  arc          = 0pt,
  outer arc    = 0pt,
  width        = 0.27\paperwidth,
  left         = 0.85em,
  right        = 0.85em,
  top          = 1.8em,
  bottom       = 1.8em,
  after        = \hspace{0.6em},
  nobeforeafter,
  enhanced,
  remember,
]
% ---- inicjały ----
\begin{center}
  \fcolorbox{goldAccent!55}{sidebar!88}{%
    \parbox{3cm}{\centering\vspace{1.1cm}
      {\fontsize{46}{46}\selectfont\bfseries\color{goldAccent} MG}\vspace{1.1cm}}
  }
\end{center}

\vspace{0.5em}
{\centering\goldline\par}
\vspace{0.5em}

% ---- kontakt ----
{\small\bfseries\color{goldAccent} \faIcon[regular]{envelope}\hspace{0.35em}KONTAKT}\\[0.2em]
{\footnotesize
\faIcon[regular]{envelope}\hspace{0.35em}\href{mailto:maciejgorski978@gmail.com}{maciejgorski978@gmail.com}\\[0.12em]
\faIcon{phone-alt}\hspace{0.35em}+48 514 299 431\\[0.12em]
\faIcon{map-marker-alt}\hspace{0.35em}Polska
}

\vspace{0.6em}
{\centering\goldline\par}
\vspace{0.6em}

% ---- online ----
{\small\bfseries\color{goldAccent} \faIcon{globe}\hspace{0.35em}ONLINE}\\[0.2em]
{\footnotesize
\faIcon{github}\hspace{0.35em}\href{https://github.com/Gorski-Maciej}{github.com/Gorski-Maciej}\\[0.12em]
\faIcon{linkedin}\hspace{0.35em}\href{https://linkedin.com/in/maciej-gorski-046176282}{LinkedIn}
}

\vspace{0.6em}
{\centering\goldline\par}
\vspace{0.6em}

% ---- języki ----
{\small\bfseries\color{goldAccent} \faIcon{language}\hspace{0.35em}JĘZYKI}\\[0.2em]
{\footnotesize
\textbf{Polski} --- ojczysty\\[0.08em]
\textbf{Angielski} --- B1/B2 (średniozaawansowany)
}

\vspace{0.6em}
{\centering\goldline\par}
\vspace{0.6em}

% ---- technologie główne ----
{\small\bfseries\color{goldAccent} \faIcon{code}\hspace{0.35em}STACK GŁÓWNY}\\[0.2em]
{\footnotesize
\setlength{\tabcolsep}{0.3em}
\begin{tabular}{ll}
\color{goldAccent}$\blacktriangleright$ & Python / FastAPI / Flask \\
\color{goldAccent}$\blacktriangleright$ & Docker / Kubernetes \\
\color{goldAccent}$\blacktriangleright$ & CI/CD (GitHub Actions) \\
\color{goldAccent}$\blacktriangleright$ & PostgreSQL / Redis / Kafka \\
\color{goldAccent}$\blacktriangleright$ & Linux / Bash \\
\color{goldAccent}$\blacktriangleright$ & React / Vite \\
\color{goldAccent}$\blacktriangleright$ & AWS \\
\end{tabular}
}

\vspace{0.6em}
{\centering\goldline\par}
\vspace{0.6em}

% ---- sieci ----
{\small\bfseries\color{goldAccent} \faIcon{network-wired}\hspace{0.35em}SIEĆ I INFRA}\\[0.2em]
{\footnotesize
\setlength{\tabcolsep}{0.3em}
\begin{tabular}{ll}
\color{goldAccent}$\blacktriangleright$ & TCP/IP, OSI, VLAN, VLSM \\
\color{goldAccent}$\blacktriangleright$ & Routing: RIP / OSPF \\
\color{goldAccent}$\blacktriangleright$ & Cisco (GNS3 / Packet Tracer) \\
\color{goldAccent}$\blacktriangleright$ & Firewall / iptables \\
\color{goldAccent}$\blacktriangleright$ & Wireshark / Monitoring \\
\color{goldAccent}$\blacktriangleright$ & CCNA (w trakcie) \\
\end{tabular}
}

\vspace{0.6em}
{\centering\goldline\par}
\vspace{0.6em}

% ---- kompetencje ----
{\small\bfseries\color{goldAccent} \faIcon{users}\hspace{0.35em}KOMPETENCJE}\\[0.2em]
{\footnotesize
\setlength{\tabcolsep}{0.3em}
\begin{tabular}{ll}
\color{goldAccent}$\blacktriangleright$ & Analityczne myślenie \\
\color{goldAccent}$\blacktriangleright$ & Samodzielność i proaktywność \\
\color{goldAccent}$\blacktriangleright$ & Szybka nauka \\
\color{goldAccent}$\blacktriangleright$ & Odpowiedzialność i dbałość \\
\color{goldAccent}$\blacktriangleright$ & Praca zespołowa \\
\color{goldAccent}$\blacktriangleright$ & Elastyczność / zaangażowanie \\
\end{tabular}
}

\vspace{0.6em}
{\centering\goldline\par}
\vspace{0.6em}

% ---- zainteresowania ----
{\small\bfseries\color{goldAccent} \faIcon{star}\hspace{0.35em}ZAINTERESOWANIA}\\[0.2em]
{\footnotesize
Sieci komputerowe i cyberbezpieczeństwo\\[0.04em]
Python i projekty open-source\\[0.04em]
IoT i automatyzacja domowa\\[0.04em]
Nowe technologie (kursy, webinaria)
}
\end{tcolorbox}
%
\begin{minipage}[t]{0.69\paperwidth}
\vspace{1.4em}
\hspace{1.4em}
\begin{minipage}[t]{0.9\textwidth}

% ======================== NAGŁÓWEK ========================
{\fontsize{34}{38}\selectfont\bfseries\color{sidebar} Maciej Górski\par}
\vspace{0.15em}
{\large\color{goldDark} Student Informatyki \textbar\ Python Developer \textbar\ Network Enthusiast\par}

\vspace{0.3em}
{\small\color{mutedText}
\faIcon[regular]{envelope}\hspace{0.3em}\href{mailto:maciejgorski978@gmail.com}{maciejgorski978@gmail.com} \hspace{1.2em}
\faIcon{phone-alt}\hspace{0.3em}+48 514 299 431 \hspace{1.2em}
\faIcon{github}\hspace{0.3em}\href{https://github.com/Gorski-Maciej}{github.com/Gorski-Maciej}
}

\vspace{0.6em}
\goldlineFull
\vspace{0.8em}

% ======================== CEL ZAWODOWY ========================
\section{\faIcon{bullseye}\hspace{0.35em}CEL ZAWODOWY}

\small
Student pierwszego roku informatyki z praktycznym doświadczeniem projektowym w budowie złożonych systemów mikrousługowych (FastAPI, Docker, Kafka, PostgreSQL, React). Łączę wiedzę z zakresu programowania w Pythonie z umiejętnościami administracji Linux, sieci komputerowych (przygotowanie do CCNA) oraz DevOps (Docker, Kubernetes, CI/CD, AWS).

Posiadam ukończone kursy z \textbf{bezpieczeństwa IT}, \textbf{IoT}, \textbf{testowania oprogramowania} i \textbf{administracji Linux}. Celem zawodowym jest rozwój jako \textbf{Python Developer}, \textbf{specjalista IT Support / DevOps} lub \textbf{administrator sieci}, oferując motywację do nauki, zaangażowanie i rzetelność w wykonywaniu powierzonych zadań.

\vspace{0.6em}
\goldlineFull
\vspace{0.7em}

% ======================== PROJEKTY ========================
\section{\faIcon{rocket}\hspace{0.35em}PROJEKTY WŁASNE}

% ---- CloudBudget ----
\cventry{2025--2026}{CloudBudget -- Platforma FinOps}{}{}{}{
\begin{minipage}[t]{\linewidth}
\small
\hangindent=0.5em
\setlength{\parskip}{0.15em}
\textbf{Stack:} FastAPI, PostgreSQL, DuckDB, RabbitMQ, Celery, React, Docker, Redis, Prometheus, Grafana\\
\textbf{Opis:} Kompletna platforma do analizy i optymalizacji kosztów chmurowych. Dashboard w React z wizualizacją wydatków, rekomendacje optymalizacyjne oparte o ML (scikit-learn), asynchroniczne przetwarzanie danych przez Celery/RabbitMQ. Testy pytest, CI/CD przez GitHub Actions, monitoring Prometheus+Grafana.
\newline{\footnotesize\color{goldDark}\href{https://github.com/Gorski-Maciej/PROJEKTS-WORK}{github.com/Gorski-Maciej/PROJEKTS-WORK}}
\end{minipage}}

% ---- NetGuardian ----
\cventry{2025--2026}{NetGuardian -- SecOps / SIEM / SOAR}{}{}{}{
\begin{minipage}[t]{\linewidth}
\small
\setlength{\parskip}{0.15em}
\textbf{Stack:} FastAPI, Kafka, TimescaleDB, DuckDB, Redis, scikit-learn, GeoIP, Docker, Prometheus, Grafana\\
\textbf{Opis:} Zaawansowany system detekcji zagrożeń sieciowych. Ingest zdarzeń przez Kafka, analiza behawioralna i regułowa, korelacja zdarzeń, detekcja anomalii (ML), enrichment GeoIP. Automatyczna reakcja przez playbooki, honeypot i SSH. Dashboard w czasie rzeczywistym z metrykami Prometheus.
\newline{\footnotesize\color{goldDark}\href{https://github.com/Gorski-Maciej/PROJEKTS-WORK}{github.com/Gorski-Maciej/PROJEKTS-WORK}}
\end{minipage}}

% ---- NetAegis ----
\cventry{2025--2026}{NetAegis -- Orkiestracja Agentów MCP}{}{}{}{
\begin{minipage}[t]{\linewidth}
\small
\setlength{\parskip}{0.15em}
\textbf{Stack:} FastAPI, MCP (Model Context Protocol), React, WebSocket, SQLite, Redis, Docker\\
\textbf{Opis:} Dwuwarstwowy system orkiestracji agentów MCP (Main + Operational) z trzema wyspecjalizowanymi agentami: NetPulse (telemetria), SecLog (logi bezpieczeństwa) i NetConfig (konfiguracja sieci). Frontend operatorski w React z komunikacją WebSocket w czasie rzeczywistym.
\newline{\footnotesize\color{goldDark}\href{https://github.com/Gorski-Maciej/PROJEKTS-WORK}{github.com/Gorski-Maciej/PROJEKTS-WORK}}
\end{minipage}}

% ---- InfraFlow ----
\cventry{2025--2026}{InfraFlow -- Monitorowanie Infrastruktury}{}{}{}{
\begin{minipage}[t]{\linewidth}
\small
\setlength{\parskip}{0.15em}
\textbf{Stack:} FastAPI, TimescaleDB, Prometheus, Grafana, Docker\\
\textbf{Opis:} Platforma do obserwowalności i automatyzacji operacyjnej infrastruktury. Metryki w czasie rzeczywistym, dashboardy Grafana, alertowanie i historyczne przechowywanie danych w TimescaleDB.
\newline{\footnotesize\color{goldDark}\href{https://github.com/Gorski-Maciej/PROJEKTS-WORK}{github.com/Gorski-Maciej/PROJEKTS-WORK}}
\end{minipage}}

\vspace{0.5em}
\goldlineFull
\vspace{0.7em}

% ======================== EDUKACJA ========================
\section{\faIcon{graduation-cap}\hspace{0.35em}EDUKACJA}

\cventry{2025--obecnie}{Informatyka (studia niestacjonarne)}{}{}{}{
\begin{minipage}[t]{\linewidth}
\small
Specjalizacja: administracja sieci i systemów
\end{minipage}}

\vspace{0.3em}
\goldlineFull
\vspace{0.7em}

% ======================== CERTYFIKATY ========================
\section{\faIcon{certificate}\hspace{0.35em}CERTYFIKATY I KURSY}

\cvitem{\small 2025}{\small\textbf{Cisco CCNA} --- w trakcie przygotowań}
\cvitem{\small 2025}{\small\textbf{Sieci komputerowe i bezpieczeństwo IT} --- kursy ukończone}
\cvitem{\small 2025}{\small\textbf{Python Developer} --- programowanie obiektowe, testy jednostkowe (pytest), Flask, FastAPI}
\cvitem{\small 2025}{\small\textbf{Linux / DevOps} --- administracja Debian/Ubuntu, Bash, Docker, CI/CD, Kubernetes}
\cvitem{\small 2025}{\small\textbf{IoT i automatyzacja} --- architektura urządzeń, integracja, podstawy}

\vspace{0.3em}
\goldlineFull
\vspace{0.7em}

% ======================== UMIEJĘTNOŚCI ========================
\section{\faIcon{cogs}\hspace{0.35em}SZCZEGÓŁOWE UMIEJĘTNOŚCI}

\cvitem{\textbf{Python}}{FastAPI, Flask, Pandas, Celery, SQLAlchemy, pytest, programowanie obiektowe, skrypty automatyzujące, analiza i wizualizacja danych}
\cvitem{\textbf{DevOps / Cloud}}{Docker, Docker Compose, Kubernetes (klastry, manifesty, sekrety), GitHub Actions (CI/CD -- build, testy, deployment), AWS (usługi chmurowe), zarządzanie rejestrami obrazów}
\cvitem{\textbf{Bazy danych}}{PostgreSQL, TimescaleDB, DuckDB, Redis, Kafka, RabbitMQ, MySQL, SQLite, modelowanie relacyjne}
\cvitem{\textbf{Monitoring}}{Prometheus, Grafana, metryki aplikacyjne, diagnostyka}
\cvitem{\textbf{Frontend}}{React 18, Vite, HTML/CSS, WebSocket, axios}
\cvitem{\textbf{Sieci}}{TCP/IP, model OSI, VLAN, routing RIP/OSPF, VLSM, Cisco (GNS3/Packet Tracer), firewall/iptables, Wireshark, monitorowanie sieci, CCNA (w trakcie)}
\cvitem{\textbf{Systemy}}{Linux (Debian/Ubuntu -- administracja, skrypty shell, cron, uprawnienia), Windows, Git}

\vspace{0.3em}
\goldlineFull
\vspace{0.7em}

% ======================== MINI-PROJEKTY ========================
\section{\faIcon{flask}\hspace{0.35em}MINI-PROJEKTY}

\cvitem{2025}{\small\textbf{HashTable implementation} --- własna implementacja struktury \texttt{key-value} z obsługą kolizji (Python, DSA)}
\cvitem{2025}{\small\textbf{Monitor sieci lokalnej} --- skrypt Bash: \texttt{ping}, \texttt{arp-scan}, logowanie wyników przez \texttt{cron}}
\cvitem{2025}{\small\textbf{Analiza logów} --- parser wyciągający statystyki błędów i adresów IP (Python, Regex)}
\cvitem{2025}{\small\textbf{Skrypt backupu} --- automatyczne archiwizowanie z rotacją kopii i logowaniem (Bash)}
\cvitem{2025}{\small\textbf{Symulacja routingu GNS3} --- topologia z OSPF, VLAN, ACL (Cisco, networking lab)}

\end{minipage}
\end{minipage}

\end{document}
